\begin{theorem}[模 \(2\) 塌缩特征多项式强制的高阶差分多项式性]\label{thm:pom-mod2-difference-binomial-basis}
令 \((a_m)_{m\ge 0}\) 为取值于 \(\FF_2\) 的序列，记前移算子 \(E\) 为 \((Ea)_m:=a_{m+1}\)，并记差分算子
\[
\Delta:=E+1
\]
（在特征 \(2\) 下与前向差分一致）。
若存在整数 \(b\ge 1\) 与阈值 \(m_0\) 使得对所有 \(m\ge m_0\) 有
\[
\Delta^b a_m=0,
\]
则存在常数 \(c_0,\dots,c_{b-1}\in\FF_2\) 使得对所有充分大的 \(m\) 都有
\[
a_m=\sum_{j=0}^{b-1} c_j\binom{m}{j}\quad(\mathrm{mod}\ 2).
\]
因此 \((a_m)\) 在 \(\FF_2\) 上最终等于一个次数 \(<b\) 的“二项式基多项式序列”。
\end{theorem}

\begin{proof}
组合恒等式给出 \(\Delta\binom{m}{j}=\binom{m}{j-1}\)（并约定 \(\binom{m}{-1}=0\)），在 \(\FF_2\) 上仍成立。
因此当且仅当 \(j<b\) 时有 \(\Delta^b\binom{m}{j}=0\)，从而任意右端线性组合必满足 \(\Delta^b a_m=0\)。
\par
反之，设对所有 \(m\ge m_0\) 有 \(\Delta^b a_m=0\)。
则 \(\Delta^{b-1}a_m\) 在 \(m\ge m_0\) 上最终常值，记该常值为 \(c_{b-1}\in\FF_2\)。
令
\(
a^{(1)}_m:=a_m-c_{b-1}\binom{m}{b-1}
\)，
则对充分大的 \(m\) 有 \(\Delta^{b-1}a^{(1)}_m=0\)。
递归下降，依次构造 \(c_{b-2},\dots,c_0\) 使得
\(
a_m-\sum_{j=0}^{b-1}c_j\binom{m}{j}
\)
的高阶差分最终为零，从而该差序列最终恒为零，得到所示表示。证毕。
\end{proof}

\begin{corollary}[碰撞矩序列的模 \(2\) 阶跃：\(q=7\) 处差分阶数从 \(2\) 跃迁到 \(4\)]\label{cor:pom-mod2-staircase-q7-difference-order-jump}
沿用注 \ref{rem:pom-charpoly-mod2-staircase-q2-10} 的记号。
在 \(q\le 10\) 的已核验范围内，最小递推特征多项式满足
\[
P_q(\lambda)\equiv \lambda^{a_q}(\lambda+1)^{b_q}\pmod 2,
\]
且 \(b_q=2\) 在 \(q=4,5,6\) 上成立，而从 \(q=7\) 起跃迁为 \(b_q=4\) 并在 \(q=7,8,9,10\) 上稳定保持。
因此对充分大的 \(m\) 有
\[
\Delta^{b_q}S_q(m+a_q)\equiv 0\pmod 2,
\]
并由定理 \ref{thm:pom-mod2-difference-binomial-basis} 得：\(S_q(m)\bmod 2\) 在 \(q=4,5,6\) 时最终为次数 \(<2\) 的二项式基多项式序列，而在 \(q=7,8,9,10\) 时最终为次数 \(<4\) 的二项式基多项式序列。
\end{corollary}

\begin{conjecture}[模 \(2\) 塌缩指数的幂次塔：二值跳跃层级无界]\label{conj:pom-mod2-collapse-exponent-tower}
设 \(b_q\) 为模 \(2\) 塌缩分解
\(
P_q(\lambda)\equiv \lambda^{a_q}(\lambda+1)^{b_q}\ (\mathrm{mod}\ 2)
\)
中 \((\lambda+1)\) 的指数（按最小整系数递推特征多项式 \(P_q\) 的模 \(2\) 约化定义）。
猜想：
\begin{enumerate}
\normalfont
  \item \(b_q\) 仅取 \(2\) 的幂并随 \(q\) 非减；
  \item 存在无穷增序列 \(q_1<q_2<\cdots\) 使得 \(b_{q_{n+1}}=2\,b_{q_n}\)；
  \item 对每个 \(q\)，\(b_q\) 给出 \(S_q(m)\bmod 2\) 的最小幂零差分阶数：对充分大的 \(m\)，有 \(\Delta^{b_q}S_q(m)=0\) 且 \(\Delta^{b_q/2}S_q(m)\not\equiv 0\)。
\end{enumerate}
该猜想把 \(q=7\) 处的指数翻倍提升为一般规律，并预言碰撞矩的二值影子携带无界的“幂零深度”。
\end{conjecture}

\endinput
